% !TEX root = ../metatime_monograph.tex
\chapter{Prospective-Observable Identifiability}
\label{app:observable_identifiability}
\label{app:prospective_observable_identifiability}


% CITATION_CONTEXT_V10_9_BEGIN
\paragraph{Established context.}
Observable selection and cancellation tests are governed by preregistration and
reproducibility principles \cite{Nosek2018,Munafo2017,Stark2018}.  Experimental
anchoring uses direct Higgs--charm constraints and current particle reference data
\cite{ATLAS2022Charm,PDG2024}; the assigned future ratios remain frozen model
predictions.
% CITATION_CONTEXT_V10_9_END

\section{Scope}

Appendix~\ref{app:up_sector_common_baseline_no_go} closed a single common
additive up-sector baseline for the fixed v10.2 trace.  The subsequent v10.5
freeze retained two possible architecture classes: a sector baseline and a
universal generation-release operator.  These two classes cannot be tested by
the same observable unless that observable is sensitive to both.

This appendix proves the required identifiability split.  It selects no new
mass formula and inspects no future Higgs-coupling likelihood.

\section{Architecture classes}

A pure sector-invariant baseline has the form

\begin{equation}
B_i=B_{s(i)},
\label{eq:sector_invariant_baseline}
\end{equation}

where all input coordinates are invariant across generations inside the
declared sector $s(i)$.  A candidate using generation, mode, generation-varying
$v_7$, or a particle-level holonomy coordinate is reclassified as a relative
release operator.

The universal relative release has the form

\begin{equation}
\Delta_f(g)=D\!\left(Z_{f,g},Z_{f,a(g)}\right),
\label{eq:universal_relative_release}
\end{equation}

with the same algebraic operator $D$ applied to charged leptons, down-type
quarks, and up-type quarks.

\section{Cancellation theorem}

The frozen v3.4 orientation layer places both charm and top in the
\path{heavy_quark_resonance} sector.  Hence a sector-invariant baseline
gives

\[
\ln y_c=B_{\rm heavy}+\Delta_c,
\qquad
\ln y_t=B_{\rm heavy}+\Delta_t.
\]

Subtracting yields

\begin{equation}
\boxed{
\ln\frac{y_c}{y_t}=\Delta_c-\Delta_t
}
\label{eq:yc_yt_baseline_cancellation}
\end{equation}

and the baseline cancels identically.  Therefore a direct charm-to-top Higgs
coupling ratio cannot identify a pure sector baseline.  It can test only a
within-sector relative-release operator.

\section{Corrected prospective mapping}

For the sector-invariant baseline class, the primary prospective observable is
the first qualifying joint ATLAS/CMS direct charm-to-tau Higgs-coupling
likelihood released after 28 July 2026.  Charm belongs to
\path{heavy_quark_resonance}, whereas tau belongs to
\path{charged_lepton_small_seed}; therefore the baseline does not cancel
from $y_c/y_\tau$.

For the universal generation-release class, the primary prospective observable
remains the first qualifying joint ATLAS/CMS direct charm-to-top coupling
likelihood released after the same date.  The frozen v10.2 structural
prediction is

\begin{equation}
\boxed{
\left|\frac{y_c}{y_t}\right|_{v10.2}
=0.006345210526463283
}
\label{eq:frozen_yc_yt_prediction}
\end{equation}

from

\[
\ln\frac{y_c}{y_t}
=7.899965155203114-12.96002015366969.
\]

Each prediction must be frozen before the corresponding likelihood is opened.
PASS requires inclusion in the published 95\% confidence region.  Failure is
retained without formula modification, and an unavailable dataset does not
license post-hoc observable replacement.

\section{Representation-source correction}

The archived v0.7 representation CSV displayed gauge-allowed transitions as
\texttt{blocked} because its residual label used an incorrect hypercharge sign.
The repository corrected this in v1.0 using

\begin{equation}
-Y_L+Y_H+Y_R=0.
\end{equation}

Candidate construction must therefore use the corrected v1.0 representation
features and not the erroneous v0.7 status column.

\section{Status}

\[
\boxed{
\text{technical PASS}
\;\land\;
\text{observable mapping corrected}
\;\land\;
\text{formula selection not begun}
}
\]

The executable audit and machine-readable ledger are located at
\path{TIR/validation/up_sector_observable_identifiability_v10_6.py}.
